% ── Dizajn a architektúra ─────────────────────────────────────
\chapter{Dizajn a architektúra}

        Táto kapitola popisuje celkovú architektúru systému, štruktúru zdrojového kódu
a~návrh používateľského rozhrania. Správny architektonický návrh je základom udržiavateľnej
a~rozšíriteľnej aplikácie.

\section{Architektúra systému}

        Aplikácia je navrhnutá ako klasická klient-server webová aplikácia s~trojvrstvovou
architektúrou:

\begin{itemize}
  \item \textbf{Prezentačná vrstva} -- HTML/CSS/JavaScript šablóny (Jinja2), ktoré
        renderuje Flask a~zasiela klientovi cez HTTPS.
  \item \textbf{Aplikačná vrstva} -- Flask blueprinty obsahujúce business logiku:
        spracovanie požiadaviek, autorizáciu, komunikáciu s~AI~modelom a~extrakciu
        textu z~dokumentov.
  \item \textbf{Dátová vrstva} -- SQLAlchemy ORM modely komunikujúce s~SQLite
        databázou. Súbory dokumentov sú uložené na filesystéme servera.
\end{itemize}

        Kód je rozdelený do samostatných Flask blueprintov podľa funkčných celkov:
\texttt{auth}, \texttt{dashboard}, \texttt{projects}, \texttt{tasks}, \texttt{documents},
\texttt{ai\_chat}, \texttt{team} a~\texttt{settings}. Každý blueprint obsahuje vlastné
routes, šablóny a~pomocné funkcie, čo umožňuje nezávislý vývoj a~testovanie jednotlivých
modulov.

\section{Štruktúra projektu}

        Zdrojový kód aplikácie je organizovaný nasledovne:

\begin{itemize}
  \item \texttt{app/} -- hlavný adresár aplikácie
  \begin{itemize}
    \item \texttt{blueprints/} -- jednotlivé moduly (auth, projects, tasks, \ldots)
    \item \texttt{models/} -- SQLAlchemy modely databázových entít
    \item \texttt{services/} -- servisná vrstva (extrakcia textu, AI~komunikácia)
    \item \texttt{templates/} -- Jinja2 HTML šablóny
    \item \texttt{static/} -- CSS, JavaScript, obrázky
  \end{itemize}
  \item \texttt{config.py} -- konfigurácia aplikácie (databáza, AI, bezpečnosť)
  \item \texttt{run.py} -- vstupný bod pre spustenie servera
  \item \texttt{requirements.txt} -- zoznam Python závislostí
  \item \texttt{uploads/} -- adresár pre nahrané dokumenty
  \item \texttt{instance/promat.db} -- SQLite databázový súbor
\end{itemize}

\begin{figure}[H]
  \centering
  % \includegraphics[width=0.7\textwidth]{images/struktura_projektu.png}
  \fbox{\parbox{0.65\textwidth}{\centering\vspace{2.5cm}
    [TBD]
  \vspace{2.5cm}}}
  \caption{Adresárová štruktúra projektu}
  \label{fig:struktura}
\end{figure}

\section{Používateľské rozhranie}

        Používateľské rozhranie je navrhnuté s~dôrazom na jednoduchosť a~prehľadnosť.
Každá stránka obsahuje fixný bočný panel (sidebar) s~navigáciou a~hlavný obsahový
priestor. Aplikácia podporuje svetlý aj tmavý režim, pričom používateľova voľba je
uložená lokálne.

        Rozhranie obsahuje nasledujúce kľúčové obrazovky:

\begin{itemize}
  \item \textbf{Dashboard} -- prehľad projektov, úloh a~aktivity
  \item \textbf{Zoznam projektov} -- prehľad všetkých projektov s~filtráciou
  \item \textbf{Detail projektu} -- úlohy, dokumenty a~členovia projektu
  \item \textbf{Kanban tabuľa} -- drag-and-drop zobrazenie úloh podľa stavu
  \item \textbf{Dokumenty} -- nahrávanie, prezeranie a~AI~chat nad dokumentmi
  \item \textbf{AI~chat} -- konverzačné rozhranie pre prácu s~dokumentmi
\end{itemize}

\begin{figure}[H]
  \centering
  % \includegraphics[width=0.9\textwidth]{images/ui_dashboard.png}
  \fbox{\parbox{0.85\textwidth}{\centering\vspace{3cm}
    [TBD]
  \vspace{3cm}}}
  \caption{Dashboard aplikácie}
  \label{fig:ui_dashboard}
\end{figure}

\begin{figure}[H]
  \centering
  % \includegraphics[width=0.9\textwidth]{images/ui_kanban.png}
  \fbox{\parbox{0.85\textwidth}{\centering\vspace{3cm}
    [TBD]
  \vspace{3cm}}}
  \caption{Kanban tabuľa projektu}
  \label{fig:ui_kanban}
\end{figure}
